\documentclass{beamer}
\usepackage[utf8]{inputenc}
\usepackage[brazilian]{babel}
\usepackage{booktabs}
\usepackage{xcolor}

% Tema CERISE — mesma identidade visual dos outros slides do projeto.
\definecolor{cerisemagenta}{HTML}{A6186B}
\definecolor{ceriserosa}{HTML}{E6186D}
\definecolor{cerisecoral}{HTML}{F5642D}
\definecolor{ceriseroxo}{HTML}{5B1A6B}

\setbeamercolor{title}{fg=cerisemagenta}
\setbeamercolor{frametitle}{fg=white,bg=cerisemagenta}
\setbeamercolor{structure}{fg=cerisemagenta}
\setbeamercolor{block title}{fg=white,bg=ceriseroxo}
\setbeamercolor{block body}{bg=cerisemagenta!5}
\setbeamercolor{alerted text}{fg=cerisecoral}
\setbeamercolor{item}{fg=cerisemagenta}

\newenvironment{numerobloco}[1]{\begin{block}{#1}}{\end{block}}
\newenvironment{decisaobloco}[1]{\begin{alertblock}{#1}}{\end{alertblock}}

\title{Progresso do Harbor}
\subtitle{Integração do dataset OpenPack: 8º dataset do chat e validação do RAG com dado real}
\author{Yan Santos Leite}
\date{13 de setembro de 2026}

\begin{document}

\frame{\titlepage}

\begin{frame}{Resumo}
  \begin{itemize}
    \item Integramos o \textbf{OpenPack} (dataset público de reconhecimento de operações de
      trabalho em embalagem logística, sensores vestíveis IMU) como o \textbf{8º dataset} do
      Harbor — o primeiro dado real usado para validar o RAG do chat com conteúdo além de
      texto puro, que até aqui só tinha sido testado contra 26 gráficos sintéticos gerados
      pelo próprio time.
    \item As três rotas do chat (\emph{contexto}, \emph{SQL}, \emph{RAG}) foram integradas e
      validadas: roteamento sem nenhuma regressão, dois gates de segurança novos, e um corpus
      de 4.000 documentos indexado a partir de 1.000 janelas de sensor.
    \item \textbf{Achado principal}: aplicamos a arquitetura \emph{RAG-HAR} (publicada em maio
      de 2026) — converter estatísticas de sensor em texto por template, sem nenhum modelo de
      visão. Ao medir esse retrieval, descobrimos que a métrica certa não é ``achar o documento
      exato'' e sim ``achar vizinhos da mesma classe'' — uma distinção formal da literatura de
      retrieval que mudou como avaliamos o resultado.
    \item RGB (fotos da operação) ficou fora do escopo por enquanto: o acesso ao dataset
      completo exige aprovação institucional externa, ainda pendente.
  \end{itemize}
\end{frame}

\begin{frame}{Por que o OpenPack, e por que agora}
  \begin{numerobloco}{A lacuna que isso fecha}
    O item 6 do plano de RAG multimodal do Harbor (a frente que trata de imagens, hoje só
    gráficos técnicos sintéticos) estava bloqueado havia dias, com a nota explícita:
    \emph{``aguardando o dataset real do usuário chegar — imagens sintéticas atuais são só
    prova de conceito''}. O OpenPack não resolve essa lacuna de imagem diretamente, mas traz
    o primeiro dado real do projeto, e sua trilha de sensores aplica a mesma filosofia
    (caption-then-embed) a um tipo de dado novo, sem exigir modelo de visão nenhum.
  \end{numerobloco}

  \pause
  \textbf{O que o OpenPack contribui:} 53,8 horas de dados reais de 16+ sujeitos realizando
  10 operações de embalagem, com sensores IMU (acelerômetro, giroscópio, quaternion) e
  anotação humana profissional — o primeiro corpus do Harbor que não foi gerado
  programaticamente pelo próprio time.

  \pause
  \begin{decisaobloco}{Contribuição concreta}
    Aplicamos o \textbf{RAG-HAR} (Sivaroopan et al., arXiv:2512.08984, maio/2026): a primeira
    arquitetura \emph{training-free} de reconhecimento de atividade humana via retrieval —
    sem nenhum modelo treinado, sem GPU. Ela converte estatísticas de sensor em texto
    estruturado e usa retrieval para dar contexto a um LLM, exatamente a filosofia que o
    Harbor já seguia em outras frentes (roteamento por regra, NL-to-SQL com self-repair).
  \end{decisaobloco}
\end{frame}

\begin{frame}{Como funciona, em uma frase}
  \small
  \textbf{Exemplo concreto:} imagine perguntar ao chat \emph{``qual operação tem o movimento
  mais intenso e irregular no punho direito?''} — o sistema não precisa de uma câmera nem de um
  modelo de visão computacional para responder. Ele lê os números do sensor vestível
  (acelerômetro, giroscópio), calcula 8 estatísticas por janela de 4 segundos (média, máximo,
  desvio-padrão, número de picos, etc.) e escreve isso como uma frase em português — depois
  busca essa frase como faria com qualquer manual técnico.

  \pause
  \begin{numerobloco}{Por que não usar um modelo de visão aqui}
    Não há imagem no IMU — a pergunta nem se coloca. Mas vale contrastar: numa frente
    \emph{diferente} do Harbor (RAG sobre gráficos técnicos, imagem de verdade), já medimos que
    \textbf{nenhum modelo de visão local disponível} atinge qualidade suficiente sem custar
    minutos por imagem (o melhor candidato, \texttt{qwen3-vl:4b}, ficou em 77,5\% de fidelidade
    e 412 segundos por imagem). Para sensores, a solução é mais simples e mais barata: o
    ``resumo em texto'' é gerado por uma fórmula matemática determinística
    (\texttt{pandas}/\texttt{numpy}), não por um modelo de IA — sem alucinação possível e sem
    custo de inferência.
  \end{numerobloco}
\end{frame}

\begin{frame}{Escala e engenharia de dados}
  \small
  \begin{numerobloco}{Números reais do pipeline}
    \begin{center}
    \begin{tabular}{lr}
      \toprule
      Sessões processadas & 102 (16 sujeitos $\times$ até 5 sessões) \\
      Sessões válidas & 102/102 (100\%) \\
      Janelas de sensor segmentadas & 78.667 \\
      Janelas amostradas para o corpus & 1.000 (100 por operação) \\
      Sujeitos representados por operação & 21 (todas as 10 operações) \\
      Documentos indexados & 4.000 (4 por janela) \\
      \bottomrule
    \end{tabular}
    \end{center}
  \end{numerobloco}

  \pause
  \textbf{Por que amostrar:} o corpus completo (78.667 janelas $\times$ 4 = mais de 300 mil
  documentos) não caberia no índice em memória nem seria viável de reavaliar a cada iteração.
  Usamos o \textbf{método do maior resto} (Hamilton/Hare-Niemeyer — o mesmo algoritmo usado
  para distribuir cadeiras em eleições proporcionais) para dividir a cota de 100 janelas por
  operação entre os 21 sujeitos sem viés de arredondamento — na primeira tentativa, uma divisão
  inteira simples havia entregado só 84 janelas por operação em vez de 100.
\end{frame}

\begin{frame}{As três rotas do chat, sem regressão}
  \small
  \begin{numerobloco}{Contexto e SQL}
    Outputs agregados do pipeline (duração por operação, transições entre operações,
    variabilidade por sujeito) alimentam a rota \emph{contexto}. Cinco tabelas novas no banco
    Postgres alimentam a rota \emph{SQL}, com uma nota explícita no schema: dados do OpenPack
    são sobre \textbf{pessoas}, nunca devem ser cruzados com tabelas de \textbf{máquina}
    (OEE, CNC, sensores industriais) — mesmo que ambos mencionem ``sensor''.
  \end{numerobloco}

  \pause
  \begin{decisaobloco}{Dois gates de segurança novos}
    Um impede o cruzamento pessoa$\times$máquina descrito acima (risco previsto antes mesmo de
    acontecer, seguindo o mesmo padrão de bugs já corrigidos em outros datasets). O outro
    recusa, de forma determinística, qualquer pedido de identificar ou avaliar um sujeito
    específico — exigência da licença do dataset (dados anonimizados por desenho), não só
    precaução.
  \end{decisaobloco}

  \pause
  \textbf{Resultado do harness de regressão:} roteamento subiu de 55/67 (82,1\%) para 58/69
  (84,1\%) — as duas perguntas novas (uma para cada gate) roteiam corretamente, e
  \textbf{nenhuma} das 67 perguntas antigas passou a falhar.
\end{frame}

\begin{frame}{O achado central: medindo a coisa certa}
  \small
  \begin{numerobloco}{O que aconteceu}
    A primeira avaliação do corpus RAG deu \textbf{Recall@5 = 0\%} — nenhuma das 45 perguntas
    com alvo conhecido encontrou a janela exata entre os 5 primeiros resultados, mesmo usando
    só busca por similaridade de texto (sem nenhum componente de palavra-chave).
  \end{numerobloco}

  \pause
  \textbf{Como investigamos:} isolamos a causa em três etapas — testamos com e sem busca por
  palavra-chave (BM25), testamos frases sem números (só nome de sujeito/sessão), e comparamos
  a pontuação bruta de discriminação de diferentes termos do texto. Todas as evidências
  apontaram para a mesma causa: como as 1.000 janelas usam o mesmo modelo de frase, com só os
  números mudando, a maioria das palavras aparece em quase todos os documentos — não sobra
  sinal de texto para diferenciar uma janela específica das outras 99 da mesma operação.

  \pause
  \begin{decisaobloco}{A correção não foi no código, foi na pergunta}
    Relendo o próprio artigo do RAG-HAR, confirmamos que ele nunca tenta achar ``o documento
    exato'' — ele busca vizinhos \emph{parecidos} para dar contexto a um classificador. A
    literatura de recuperação de informação chama isso de \textbf{retrieval por categoria}
    (achar algo da mesma classe), diferente de \textbf{retrieval por instância} (achar o item
    exato) — nosso teste original media a pergunta errada para esta arquitetura.
  \end{decisaobloco}
\end{frame}

\begin{frame}{Resultado final, com a métrica certa}
  \small
  \begin{numerobloco}{k-NN label purity: fração dos vizinhos que são da mesma operação}
    \begin{center}
    \begin{tabular}{lr}
      \toprule
      Pureza média (5 vizinhos, 45 perguntas) & 12,9\% \\
      Acaso esperado (10 operações balanceadas) & 10,0\% \\
      Perguntas com pureza perfeita (100\%) & 1/45 \\
      Perguntas com pureza zero (0\%) & 33/45 \\
      \bottomrule
    \end{tabular}
    \end{center}
  \end{numerobloco}

  \pause
  \textbf{O que isso significa:} o retrieval discrimina \emph{um pouco} acima do acaso, não
  fortemente. É um resultado honesto, não um sucesso nem um fracasso — mostra que estatísticas
  de primeira ordem (média, desvio, número de picos) sobre 6 canais de movimento têm poder
  discriminativo real, mas fraco, entre 10 operações de embalagem parecidas entre si. Uma
  comparação direta e justa com o baseline anterior do RAG multimodal (26 imagens sintéticas,
  Recall@3 = 38\%) não é possível — as métricas medem coisas diferentes por desenho —, mas o
  processo de chegar a essa conclusão é, em si, um resultado metodológico que vale registrar.
\end{frame}

\begin{frame}{Riscos e pendências}
  \small
  \begin{itemize}
    \item[\textcolor{green!60!black}{\textbullet}] \textbf{Rotas contexto/SQL} — fechadas,
      validadas sem regressão no harness completo.
    \item[\textcolor{green!60!black}{\textbullet}] \textbf{Corpus RAG e métrica de avaliação}
      — fechados, com achado metodológico documentado.
    \item[\textcolor{orange!90!black}{\textbullet}] \textbf{Poder discriminativo do retrieval
      (12,9\%)} — funcional, mas fraco. Não bloqueia a integração; abre uma pergunta em aberto
      sobre se features de ordem mais alta ajudariam, ainda sem teste que confirme.
    \item[\textcolor{red!80!black}{\textbullet}] \textbf{Imagens RGB (Fase 5)} — bloqueada por
      aprovação institucional externa, ainda não solicitada. O único frame disponível sem essa
      aprovação cobre um instante antes do início de qualquer operação anotada, insuficiente
      para o teste original planejado (12 imagens, uma por operação). O mecanismo técnico foi
      validado apesar disso: um \emph{prompt} de legendagem herdado (feito para gráficos, não
      fotos) gerava legenda vazia; corrigido, o mesmo modelo descreveu a cena real
      corretamente.
    \item[\textcolor{orange!90!black}{\textbullet}] \textbf{Benchmark acadêmico comparável}
      (meta BRACIS 2027) — os números desta sessão são a base, mas a comparação formal contra
      os resultados publicados de HAR no artigo original do OpenPack ainda não foi feita.
  \end{itemize}
\end{frame}

\begin{frame}{Próximos passos}
  \begin{enumerate}
    \item \textbf{Solicitar acesso ao RGB completo} — preencher o formulário institucional dos
      autores do OpenPack (pendente, é uma ação pessoal, não automatizável). Critério de
      conclusão: aprovação recebida e ao menos 1 keyframe por operação disponível.
    \item \textbf{Investigar o poder discriminativo fraco} — é uma hipótese em aberto, não
      testada nesta sessão, se features de ordem mais alta (ex.\ frequência dominante via FFT)
      melhorariam a pureza acima de 12,9\%; vale medir antes de decidir se a abordagem atual
      já chegou ao teto.
    \item \textbf{Comparação formal com o benchmark do OpenPack} — reproduzir/contrastar contra
      os números de reconhecimento de atividade já publicados no artigo original, usando a
      mesma filosofia \emph{training-free} do RAG-HAR.
    \item \textbf{Sincronizar o dashboard} — expor a nova aba de chat do OpenPack em
      \texttt{dashboard/app.py}, seguindo a mesma disciplina já aplicada aos outros 4 datasets
      ativos.
  \end{enumerate}
\end{frame}

\end{document}
